%% ma: L11-B16-0062
%% lop: 11
%% chuong: 5
%% bai: 16
%% dang: 11.16.4
%% loai: DS
%% muc: [TH, TH, VD, VD]
%% dap_an: "ĐSSS"
%% nguon: "tai-lieu-hoc-tap-mon-toan-lop-11-theo-chuong-trinh-giao-duc-pho-thong-2018.pdf (Lớp Toán Thầy Duy, Phạm Lê Duy, THPT Chu Văn An), Chương 3 Bài 2, Phần 2, Câu 6 (tr.378)"
%% kien_thuc: "Lớp 11 Bài 16 (giới hạn của hàm số)"
%% trang_thai: chinh-thuc
%% ghi_chu: "Bảng đáp án gốc: a Đ, b S, c Đ, d Đ. Theo L = (a-1)/b thì c), d) là Sai (kho giải ra Đ S S S). Có thể đề gốc in nhầm dấu ở c), d) (a - 2b = 1, a - b = 1 thì Đúng); nhờ giáo viên đối chiếu bản gốc. Đã duyệt 03/11/2026."
\begin{ex}
Cho $a,b$ là các số thực khác $0$. Biết giới hạn $L=\lim\limits_{x\to-\infty}\dfrac{\sqrt{x^2-3x}+ax}{bx-1}$. Khi đó
\choiceTF
{\True $L=3$ khi $\dfrac{a-1}b=3$.}
{$L=6$ khi $\dfrac{a-1}b=4$.}
{$L=2$ khi $a+2b=1$.}
{$L=1$ khi $a+b=1$.}
\loigiai{\sloppy Với $x<0$: $\sqrt{x^2-3x}=-x\sqrt{1-\frac3x}$. Chia tử và mẫu cho $x$: $L=\lim\dfrac{-\sqrt{1-\frac3x}+a}{b-\frac1x}=\dfrac{a-1}b$.
a) $L=3\Leftrightarrow\dfrac{a-1}b=3$. Đúng.
b) $L=6\Leftrightarrow\dfrac{a-1}b=6\ne4$. Sai.
c) $L=2\Leftrightarrow a-1=2b\Leftrightarrow a-2b=1$; với $a+2b=1$ thì $L=\dfrac{-2b}b=-2\ne2$. Sai.
d) $L=1\Leftrightarrow a-b=1$; với $a+b=1$ thì $L=\dfrac{-b}b=-1\ne1$. Sai.}
\end{ex}
